\begin{Theorem}{Euler-Lagrange equations}
\label{thm:elasticity-euler-lagrange-equations}
Suppose $\overline{\bm{p}}\in\mathcal{D}_{\mathcal{E}}\cap C^2(\overline\Omega)$ solves \eqref{eq:elasticity-equilibrium-global-minimization} and that \eqref{eq:elasticity-positive-orientation-determinant} holds almost everywhere in $\Omega$. Then at every $\bm{x}\in\Omega$, $\overline{\bm{p}}$ must satisfy the Euler-Lagrange equations
\begin{equation}\label{eq:elasticity-euler-lagrange-equations}
-\sum_j\dfrac{\partial}{\partial x_j}\mathcal{A}_{ij}(\bm{x})+\dfrac{\partial\Psi}{\partial p_i}(\bm{p}(\bm{x}))=0,\quad i=1,2,3,
\end{equation}
where
\begin{equation}\label{eq:elasticity-euler-lagrange-stress-coefficients}
\mathcal{A}_{ij}(\bm{x}):=\dfrac{\partial\widehat{\mathcal{W}}}{\partial F_{ij}}(\bm{x},\bm{F}(\bm{x})).
\end{equation}
In a homogeneous material: $\mathcal{W}=\widehat{\mathcal{W}}(\bm{F})$, we have
\begin{equation}\label{eq:elasticity-homogeneous-stress-divergence}
\sum_j\dfrac{\partial}{\partial x_j}\mathcal{A}_{ij}(\bm{x})=\sum_{j,k,l}a_{ijkl}(\bm{x})\dfrac{\partial^2p_k}{\partial x_j\partial x_l},
\end{equation}
where
\begin{equation}\label{eq:elasticity-homogeneous-second-derivative-coefficients}
a_{ijkl}(\bm{x}):=\dfrac{\partial^2\mathcal{W}}{\partial F_{ij}\partial F_{kl}}(\bm{F}(\bm{x})).
\end{equation}

\end{Theorem}

\begingroup
\renewcommand{\endProof}{\par\nobreak\hfill$\blacksquare$\par\vspace{0.1cm}}
\begin{Proof}
Since $\overline{\bm{p}}\in\mathcal{D}_{\mathcal{E}}\cap C^2(\overline\Omega)$, we see that for any $\bm{\phi}=\sum\phi_i\bm{e}_i\in[\mathcal{D}(\Omega)]^3$ we have
\begin{equation}\label{eq:elasticity-compact-variation-admissibility}
\overline{\bm{p}}+\epsilon\bm{\phi}\in\mathcal{D}_{\mathcal{E}}
\end{equation}
for $\epsilon$ sufficiently small. Thus, for any $\bm{\phi}$, the real-valued function
\begin{equation}\label{eq:elasticity-scalar-variation-energy}
f(\epsilon):=\mathcal{E}(\overline{\bm{p}}+\epsilon\bm{\phi})
\end{equation}
is well defined in an interval about $\epsilon=0$. Furthermore, $f$ is minimized at $\epsilon=0$ since by hypothesis $\mathcal{E}$ is minimized at $\overline{\bm{p}}$. Now, using standard results on uniform convergence to take the derivative under the integral, we see that $f$ is differentiable. We get
\begin{equation}\label{eq:elasticity-first-variation-integral}
\begin{aligned}
 f^{\prime}(\epsilon)=\int_\Omega&\sum_{i,j}\dfrac{\partial\mathcal{W}}{\partial F_{ij}}(\bm{x},\overline{\bm{F}}(\bm{x})+\epsilon\nabla\bm{\phi}(\bm{x}))\dfrac{\partial\phi_i}{\partial x_j}\\
 &+\sum_i\dfrac{\partial\Psi}{\partial p_i}(\overline{\bm{p}}+\epsilon\bm{\phi})\phi_i\,\mathrm{d}\bm{x}.
 \end{aligned}
\end{equation}
Setting the derivative equal to zero at $\epsilon=0$ and integrating by parts gives us
\begin{equation}\label{eq:elasticity-first-variation-integrated-by-parts}
\int_\Omega\sum_i\left[-\sum_j\dfrac{\partial}{\partial x_j}\dfrac{\partial\mathcal{W}}{\partial F_{ij}}(\bm{x},\overline{\bm{F}}(\bm{x}))+\dfrac{\partial\Psi}{\partial p_i}(\overline{\bm{p}}(\bm{x}))\right]\phi_i(\bm{x})\,\mathrm{d}\bm{x}=0.
\end{equation}
Since $\bm{\phi}$ is arbitrary, this implies that \eqref{eq:elasticity-euler-lagrange-equations} is satisfied in the sense of distributions. Since each term in the equation is continuous, it must, in fact, be satisfied pointwise.
\end{Proof}
\endgroup
